%----------------------------------------------------------------------------------------
%    PACKAGES AND THEME
%----------------------------------------------------------------------------------------

\documentclass[aspectratio=169,xcolor=dvipsnames]{beamer}
\usetheme{SimplePlus}

\usepackage[T1]{fontenc}
\usepackage{lmodern}
\usepackage{hyperref}
\usepackage{graphicx}
\usepackage{booktabs}
\usepackage{amsmath}
\usepackage{amssymb}
\usepackage{listings}

\graphicspath{
    {./figures/}
    {./figures/schedule/}
}

\definecolor{mznKw}{RGB}{0,90,160}
\definecolor{mznCom}{RGB}{120,120,120}
\definecolor{mznStr}{RGB}{170,55,55}

\lstdefinelanguage{minizinc}{
  morekeywords={
    include, par, var, array, of, int, bool, set, constraint,
    forall, exists, sum, alldifferent, solve, satisfy, minimize, maximize,
    if, then, else, endif, in, div, mod, where, output, predicate, function
  },
  sensitive=true,
  morecomment=[l]{\%},
  morestring=[b]"
}

\lstset{
  language=minizinc,
  basicstyle=\ttfamily\scriptsize,
  keywordstyle=\color{mznKw}\bfseries,
  commentstyle=\color{mznCom}\itshape,
  stringstyle=\color{mznStr},
  breaklines=true,
  showstringspaces=false,
  frame=single,
  framesep=3pt,
  rulecolor=\color{gray!40}
}

\setbeamertemplate{footline}[frame number]

%----------------------------------------------------------------------------------------
%    TITLE PAGE
%----------------------------------------------------------------------------------------

\title{Robot Movement Planning}
\subtitle{Constraint Programming}
\author{Alessio Russo}

\institute{
    Department of Mathematical, Physical and Computer Sciences\\
    University of Parma
}
\date{}

%----------------------------------------------------------------------------------------
%    PRESENTATION SLIDES
%----------------------------------------------------------------------------------------

\begin{document}

\begin{frame}[plain,noframenumbering]
    \titlepage
\end{frame}

%------------------------------------------------

\begin{frame}{Model Formulation}
    \small
    The robot movement problem is modeled as a feasibility problem.
    For each pallet layer, the model decides how the incoming boxes are picked
    and how the resulting groups are released on the pallet.

    \begin{columns}[T,totalwidth=\textwidth]
        \begin{column}{0.50\textwidth}
            \begin{block}{Decision variables}
                \begin{itemize}
                    \item which consecutive boxes form each pick group;
                    \item when each picked group is released;
                    \item which gripper plate opens at release.
                \end{itemize}
            \end{block}
        \end{column}

        \begin{column}{0.46\textwidth}
            \begin{block}{Constraints}
                \begin{itemize}
                    \item groups have no gaps or overlaps;
                    \item each box is picked and placed exactly once;
                    \item each pick fits within the gripper size;
                    \item each group stays within one pallet row;
                    \item the opening side is free at release.
                \end{itemize}
            \end{block}
        \end{column}
    \end{columns}

    \vspace{0.25cm}
    \begin{alertblock}{Goal}
        Find any assignment that defines a feasible pick-and-release plan.
    \end{alertblock}
\end{frame}

%------------------------------------------------

\begin{frame}{Pick and Release Operation}
    \begin{figure}
        \centering
        \includegraphics[width=0.95\textwidth,height=0.78\textheight,keepaspectratio]{pick_and_release.pdf}
    \end{figure}
\end{frame}

%------------------------------------------------

\begin{frame}{Grouping Model Construction}
    \small
    The inputs of the model are the pallet grid, the box size, and the gripper capacity.
    From this data, the model derives the number of boxes that can be picked at once
    and the number of groups to schedule.

    \vspace{0.20cm}
    \begin{columns}[T,totalwidth=\textwidth]
        \begin{column}{0.48\textwidth}
            \begin{block}{Input data}
                \begin{itemize}
                    \item pallet grid dimensions;
                    \item box length, used for gripper capacity;
                    \item gripper size.
                \end{itemize}
            \end{block}
        \end{column}

        \begin{column}{0.48\textwidth}
            \begin{block}{Derived quantities}
                \begin{itemize}
                    \item maximum boxes per pick;
                    \item groups needed per row;
                    \item total number of groups.
                \end{itemize}
            \end{block}
        \end{column}
    \end{columns}

    \vspace{0.30cm}
    These derived quantities define the finite domains used later in the model.
\end{frame}

%------------------------------------------------

\begin{frame}[fragile]{Model Parameters}

    \vspace{0.15cm}
    \begin{lstlisting}
include "globals.mzn";

par int: n_boxes;
par int: boxes_x_dim;
par int: boxes_along_x;
par int: boxes_along_y;
par int: grip_size;

par int: max_pkble_boxes =
  (grip_size + boxes_x_dim) div boxes_x_dim;

par int: groups_per_row =
  (boxes_along_x + max_pkble_boxes - 1) div max_pkble_boxes;

par int: n_groups = boxes_along_y * groups_per_row;
    \end{lstlisting}
\end{frame}

%------------------------------------------------

\begin{frame}{Possible Group Representation}
    \small
    The chosen representation should be compact, while preserving the fact that a pick is a consecutive interval
    in the belt order.

    \vspace{0.25cm}
    \begin{tabular}{p{0.26\textwidth} p{0.36\textwidth} p{0.30\textwidth}}
        \toprule
        \textbf{Encoding} & \textbf{Representation} & \textbf{Limitation} \\
        \midrule
        Enumerated groups
        & Store, for each group, the list of boxes it contains, using padding for shorter groups.
        & Group structure is explicit, but many variables and extra constraints are needed. \\
\midrule
        Binary assignment
        & Use one Boolean variable for each pair \((g,i)\), equal to true if box \(i\) is assigned to group \(g\).
        & Membership is explicit, but the encoding is large; contiguity and row consistency require extra constraints. \\
\midrule
        Interval encoding
        & Store only the first box and the length of each group.
        & \textbf{Adopted}: compact, and contiguity is implicit. \\
        \bottomrule
    \end{tabular}
\end{frame}


%------------------------------------------------

\begin{frame}[fragile]{Grouping Implementation}
    \vspace{0.10cm}
    \textbf{Variables:} define the interval associated with each group.
    \begin{lstlisting}
array[1..n_groups] of var 1..n_boxes: group_start;
array[1..n_groups] of var 1..max_pkble_boxes: group_len;
    \end{lstlisting}

    \vspace{0.10cm}
    \textbf{Start condition:} the first group starts from the first box.
    \begin{lstlisting}
constraint group_start[1] = 1;
    \end{lstlisting}

    \vspace{0.10cm}
    \textbf{No gaps or overlaps:} each group starts immediately after the previous one.
    \begin{lstlisting}
constraint forall(g in 1..n_groups - 1)(
  group_start[g + 1] = group_start[g] + group_len[g]
);
    \end{lstlisting}

    \vspace{0.10cm}
    \textbf{Coverage:} the last group must end at the last box.
    \begin{lstlisting}
constraint group_start[n_groups] + group_len[n_groups] - 1 = n_boxes;
    \end{lstlisting}
\end{frame}

%------------------------------------------------

\begin{frame}[fragile]{Same-Row Constraint}
    \small
    The domain of \(\texttt{group\_len}\) already enforces the gripper capacity.
    The remaining geometric constraint is that each picked interval must stay inside one pallet row.

    \vspace{0.15cm}
    \textbf{Row index of the first box:}
    \[
        (\texttt{group\_start[g]} - 1) \; \texttt{div} \; \texttt{boxes\_along\_x}
    \]

    \vspace{0.05cm}
    \textbf{Row index of the last box:}
    \[
        (\texttt{group\_start[g]} + \texttt{group\_len[g]} - 2) \; \texttt{div} \; \texttt{boxes\_along\_x}
    \]

    \vspace{0.10cm}
    \textbf{MiniZinc implementation:} the two row indices must be equal.
    \begin{lstlisting}
constraint forall(g in 1..n_groups)(
  (group_start[g] - 1) div boxes_along_x =
  (group_start[g] + group_len[g] - 2) div boxes_along_x
);
    \end{lstlisting}
\end{frame}

%------------------------------------------------

\begin{frame}{Grouping Example}
    \small
    \begin{columns}[c,totalwidth=\textwidth]
        \begin{column}{0.60\textwidth}
            \centering
            \includegraphics[width=\linewidth,height=0.72\textheight,keepaspectratio]{sample_pallet.pdf}
        \end{column}
        \begin{column}{0.36\textwidth}
            \begin{block}{Example instance}
                \(70 \times 70\) mm boxes on an \(800 \times 400\) mm pallet give an \(11 \times 5\) grid, for \(55\) target positions.
            \end{block}
        \end{column}
    \end{columns}
\end{frame}

%------------------------------------------------

\begin{frame}[fragile]{Scheduling Variables}
    \small
    Once groups are fixed, the model assigns each group to one release step.
    It also selects which plate opens, represented as a binary decision.

    \vspace{0.15cm}
    \begin{lstlisting}
array[1..n_groups] of var 1..n_groups: drop_order;
array[1..n_groups] of var 0..1: plate;

constraint alldifferent(drop_order);
    \end{lstlisting}

    \vspace{0.25cm}
    The \(\texttt{alldifferent}\) constraint makes \(\texttt{drop\_order}\) a permutation of the release steps.
    The value \(\texttt{plate[g]}\) selects which side must be checked for free space when group \(g\) is released.
\end{frame}

%------------------------------------------------

\begin{frame}[fragile]{Release Constraints: Internal Rows}
    \small
    For internal rows, both vertical neighbors exist.
    If the upper side is selected, the upper neighbor must be scheduled later; if the lower side is selected, the lower neighbor must be scheduled later.

    \vspace{0.15cm}
    \begin{lstlisting}
constraint forall(g in 1..n_groups where
  g > groups_per_row /\ g <= n_groups - groups_per_row
)(
  (
    plate[g] = 0 /\
    drop_order[g - groups_per_row] > drop_order[g]
  )
  \/
  (
    plate[g] = 1 /\
    drop_order[g + groups_per_row] > drop_order[g]
  )
);
    \end{lstlisting}
\end{frame}

%------------------------------------------------

\begin{frame}[fragile]{Release Constraints: Boundary Rows}
    \small
    The first and last pallet rows are simpler because one side is the outside of the pallet.
    Separating boundary cases also avoids invalid array accesses.

    \vspace{0.15cm}
    \begin{lstlisting}[basicstyle=\ttfamily\tiny]
constraint forall(g in 1..n_groups where g <= groups_per_row)(
  plate[g] = 0 \/
  (
    plate[g] = 1 /\
    drop_order[g + groups_per_row] > drop_order[g]
  )
);

constraint forall(g in 1..n_groups where g > n_groups - groups_per_row)(
  (
    plate[g] = 0 /\
    drop_order[g - groups_per_row] > drop_order[g]
  )
  \/
  plate[g] = 1
);
    \end{lstlisting}
\end{frame}

%------------------------------------------------

\begin{frame}{Scheduling Example}
    \centering
    \setlength{\tabcolsep}{2pt}
    \renewcommand{\arraystretch}{0.85}

    \begin{tabular}{ccc}
        \includegraphics[width=0.33\textwidth,height=0.32\textheight,keepaspectratio]{schedule_step_01.pdf} &
        \includegraphics[width=0.33\textwidth,height=0.32\textheight,keepaspectratio]{schedule_step_04.pdf} &
        \includegraphics[width=0.33\textwidth,height=0.32\textheight,keepaspectratio]{schedule_step_07.pdf} \\
        \scriptsize Step 1 & \scriptsize Step 4 & \scriptsize Step 7 \\[-0.05cm]

        \includegraphics[width=0.33\textwidth,height=0.32\textheight,keepaspectratio]{schedule_step_10.pdf} &
        \includegraphics[width=0.33\textwidth,height=0.32\textheight,keepaspectratio]{schedule_step_13.pdf} &
        \includegraphics[width=0.33\textwidth,height=0.32\textheight,keepaspectratio]{schedule_step_15.pdf} \\
        \scriptsize Step 10 & \scriptsize Step 13 & \scriptsize Step 15
    \end{tabular}

    \vspace{0.10cm}
    \small
    The red border marks the group being released; the arrow indicates the opening side.
\end{frame}

%------------------------------------------------

\begin{frame}[fragile]{Solving from Python}
    \small
    The model is solved as a feasibility problem.  Python builds each benchmark instance, sends the parameters to MiniZinc, and reads the four decision arrays from the result.

    \vspace{0.15cm}
    \begin{columns}[T,totalwidth=\textwidth]
        \begin{column}{0.32\textwidth}
            \begin{lstlisting}
solve satisfy;
            \end{lstlisting}

            \vspace{0.2cm}
            No objective function is needed: any satisfying assignment is a valid pick-and-release plan.
        \end{column}
        \begin{column}{0.63\textwidth}
            \begin{lstlisting}[language=Python,basicstyle=\ttfamily\tiny]
model    = Model("model.mzn")
solver   = Solver.lookup("gecode")
instance = Instance(solver, model)

instance["n_boxes"]       = pallet_info["n_boxes"]
instance["boxes_x_dim"]   = boxes_x_dim
instance["boxes_along_x"] = pallet_info["boxes_along_x"]
instance["boxes_along_y"] = pallet_info["boxes_along_y"]
instance["grip_size"]     = grip_size

result = instance.solve()
            \end{lstlisting}
        \end{column}
    \end{columns}
\end{frame}

%------------------------------------------------

\begin{frame}{Benchmark Setup}
    \small
    The benchmark tests the same model across several pallet geometries, box sizes, gripper sizes, and solver backends.

    \vspace{0.25cm}
    \begin{columns}[T,totalwidth=\textwidth]
        \begin{column}{0.48\textwidth}
            \begin{table}
                \centering
                \small
                \begin{tabular}{lr}
                    \toprule
                    \textbf{Pallet} & \textbf{Dimensions (mm)} \\
                    \midrule
                    Europallet & \(1200 \times 800\) \\
                    Industrie & \(1200 \times 1000\) \\
                    US 40\(\times\)48 & \(1219 \times 1016\) \\
                    Half-pallet & \(800 \times 600\) \\
                    Australian & \(1165 \times 1165\) \\
                    \bottomrule
                \end{tabular}
            \end{table}
        \end{column}
        \begin{column}{0.48\textwidth}
            \begin{itemize}
                \item Box sizes: \(50 \times 50\) to \(200 \times 200\) mm.
                \item Solvers: \texttt{gecode}, \texttt{cp-sat}, \texttt{chuffed}, \texttt{highs}.
                \item Timeout: \(300\) seconds.
                \item Gripper sizes: \(S \in \{1,2,3\}\cdot\texttt{boxes\_x\_dim}\).
            \end{itemize}
        \end{column}
    \end{columns}
\end{frame}

%------------------------------------------------

\begin{frame}{Scaling by Gripper Size}
    \small
    \begin{columns}[T,totalwidth=\textwidth]
        \begin{column}{0.36\textwidth}
            Larger grippers reduce the number of groups that must be scheduled.
            In the benchmark, the average number of groups drops from about \(75\) for \(k=1\) to about \(40\) for \(k=3\).

            \vspace{0.25cm}
            This is reflected in the runtime curves: \texttt{gecode} remains the most stable backend, while \texttt{highs} is consistently slower.
        \end{column}
        \begin{column}{0.60\textwidth}
            \begin{figure}
                \centering
                \includegraphics[width=\linewidth,height=0.68\textheight,keepaspectratio]{scaling_by_grip.pdf}
            \end{figure}
        \end{column}
    \end{columns}
\end{frame}

%------------------------------------------------

\begin{frame}{Overall Solver Behaviour}
    \small
    \begin{columns}[T,totalwidth=\textwidth]
        \begin{column}{0.40\textwidth}
            The aggregate plot confirms that propagation-based solvers are the best fit for this model.
            \texttt{gecode} stays near the lowest runtime range across almost all instances.

            \vspace{0.25cm}
            \texttt{chuffed} is competitive but less regular, while \texttt{cp-sat} shows a stronger dependence on instance size.
            Several \texttt{highs} runs approach or reach timeout.
        \end{column}
        \begin{column}{0.56\textwidth}
            \begin{figure}
                \centering
                \includegraphics[width=\linewidth,height=0.68\textheight,keepaspectratio]{scaling_overall_n_boxes.pdf}
            \end{figure}
        \end{column}
    \end{columns}
\end{frame}

%------------------------------------------------

\begin{frame}{Flattening and Execution Overhead}
    \small
    \begin{columns}[T,totalwidth=\textwidth]
        \begin{column}{0.40\textwidth}
            The total time measured by Python includes more than solver search.
            MiniZinc must flatten the model into FlatZinc, and the benchmark also pays for instance construction and process startup.

            \vspace{0.25cm}
            This overhead is most visible on very fast runs: when the solver time is tiny, even a small fixed cost separates the point from the diagonal.
        \end{column}
        \begin{column}{0.56\textwidth}
            \begin{figure}
                \centering
                \includegraphics[width=\linewidth,height=0.68\textheight,keepaspectratio]{wall_vs_solve.pdf}
            \end{figure}
        \end{column}
    \end{columns}
\end{frame}

%------------------------------------------------

\begin{frame}{Conclusions}
    \small
    \begin{itemize}
        \item The movement planning problem is naturally split into grouping and release scheduling.
        \item The interval encoding keeps grouping compact while preserving the geometric structure of the pallet rows.
        \item Release feasibility can be expressed with local constraints on vertical neighbors and release order.
        \item The benchmark favors CP-oriented solvers, especially \texttt{gecode}.
        \item Larger grippers reduce the number of groups and usually improve runtime.
    \end{itemize}

    \vspace{0.3cm}
    \begin{block}{Final point}
        Any satisfying assignment gives a complete pick-and-release plan that places every box once and releases every group toward free space.
    \end{block}
\end{frame}

%------------------------------------------------

\begin{frame}[plain,noframenumbering]
    \centering
    \Large{\textbf{Thank you for your attention!}}

    \vspace{1.2cm}
    \footnotesize
    Alessio Russo\\
    University of Parma
\end{frame}

%----------------------------------------------------------------------------------------

\end{document}
